% ---------------------------------------------------------------------------
% Banca 1 — deck Beamer (tema PMME)
%
% Identidade visual: docs/07_apresentacoes/banca1/deck_beamer/beamerthemePMME.sty
% (paleta, navbar, frametitle, rodape e utilitarios \fonte, \lead, \num,
% destaque). Capa e sumario: pmmecapa.sty e pmmesumario.sty, no mesmo diretorio.
%
% Artefato DERIVADO de docs/07_apresentacoes/banca1/02_conteudo_slides.md.
% Divergencia entre este deck e o documento de conteudo e erro do deck.
%
% Compilar SEMPRE a partir da raiz do repositorio:
%     bash scripts/apresentacao/build_deck_beamer.sh
% ---------------------------------------------------------------------------
\documentclass[aspectratio=169,11pt]{beamer}

\usepackage[T1]{fontenc}
\usepackage[utf8]{inputenc}
\usepackage[brazil]{babel}
\usepackage{lmodern}
\usepackage{amsmath}
\usepackage{amssymb}
\usepackage{booktabs}
\usepackage{graphicx}
\usepackage{array}
\usepackage{tabularx}
\usepackage{makecell}
\usepackage{ragged2e}

% --- figuras: caminhos relativos a raiz do repositorio ----------------------
\graphicspath{{output/apresentacao_banca1/}}

% --- tema -------------------------------------------------------------------
% Tudo o que e aparencia (cores, navbar, frametitle, rodape, \fonte, \lead,
% \num, destaque, colunas L/Z, \arraystretch, caption) mora no tema.
\usetheme{PMME}

% --- metadados --------------------------------------------------------------
\title{Desvantagens territoriais na escolha locacional de médicos especialistas}
\subtitle{Um modelo microeconômico para o preenchimento de vagas do Programa
Mais Médicos Especialistas}
% Slots da capa (documento canonico, slide 1). Preencher com os nomes reais
% antes da banca; o deck compila com os placeholders.
\author{Autoria}
\orientacao{Orientador(a) \textperiodcentered\ coorientador(a)}
\institute{Insper}
\date{Banca 1 \textperiodcentered\ data da banca}
% Logo e foto da capa: caminhos relativos a raiz do repositorio (o build
% compila a partir dela). Ausentes, a capa usa os substitutos do pacote.
\pmmelogo{docs/07_apresentacoes/banca1/deck_beamer/assets/logo_insper.png}
\pmmefoto{docs/07_apresentacoes/banca1/deck_beamer/assets/capa_foto.jpg}

% --- sumario: as seis secoes e a descricao curta de cada uma ---------------
% Conteudo canonico: 02_conteudo_slides.md, slide 2.
\pmmesumariofundo{docs/07_apresentacoes/banca1/deck_beamer/assets/sumario_fundo.jpg}
\pmmesumarioitem{Motivação}{Por que a pergunta importa}
\pmmesumarioitem{Pergunta}{A pergunta de pesquisa}
\pmmesumarioitem{Literatura teórica}{As três referências que fundamentam o modelo}
\pmmesumarioitem{Modelo microeconômico}{A decisão do médico e o custo do lugar}
\pmmesumarioitem{Hipótese}{O que o modelo prevê}
\pmmesumarioitem{Viabilidade empírica}{O que os dados permitem testar}

\begin{document}

% ===========================================================================
% 1 — Capa
% ===========================================================================
\rastreio{}
\begin{frame}[plain]
  \pmmecapa
\end{frame}

% ===========================================================================
% 2 — Sumário
% ===========================================================================
\rastreio{}
\begin{frame}[plain]
  \pmmesumario
\end{frame}

% ===========================================================================
% 3 — Especialistas faltam no interior, e o médico sabe por quê   (1 frame)
% ===========================================================================
\section{Motivação}
\subsection{O problema}
\rastreio{1. Motivação \textperiodcentered\ O problema}

\begin{frame}{Especialistas faltam no interior, e o médico sabe por quê}
  \only<1>{%
    \vspace{0.1em}
    {\small O Brasil tinha, em 2024, \num{353 mil médicos especialistas} --- 59\%
    dos 597 mil médicos do país. \num{55\%} estão no \textbf{Sudeste} e
    \num{6\%} no \textbf{Norte}; são \num{453} por 100 mil habitantes no
    \textbf{Distrito Federal} e \num{68} no \textbf{Maranhão}. \lead{O problema
    não é o número. É onde eles estão.}}

    \vspace{0.35em}
    {\small \lead{E o problema tem outro lado: a decisão do médico.} Para ele,
    ``município vulnerável'' não é índice --- são \textbf{desvantagens
    concretas}, e \lead{duas nós medimos}.}

    \vspace{0.3em}
    {\footnotesize\renewcommand{\arraystretch}{1.0}
    \begin{tabularx}{\textwidth}{@{}>{\bfseries\raggedright\arraybackslash}p{0.245\textwidth} L >{\raggedright\arraybackslash}p{0.215\textwidth}@{}}
      \toprule
      \normalfont\bfseries Desvantagem & \bfseries O que significa para o médico
        & \bfseries Medimos? \\
      \midrule
      Retaguarda profissional & sem segunda opinião, sem plantão, sem a quem
        encaminhar & \textbf{sim} --- CNES \\
      Infraestrutura & equipamento, insumos e equipe escassos: atender cansa
        mais e resolve menos & \textbf{em parte} --- CNES \\
      Distância da família & viver longe da família e de onde se formou &
        \textbf{não} \\
      Mercado privado ausente & a remuneração se reduz ao que o programa paga &
        \textbf{não} \\
      \bottomrule
    \end{tabularx}}
  }%
  \only<2>{%
    \vspace{0.8em}
    \begin{block}{O peso de cada uma, na literatura}
      \small
      \textbf{Brasil}, 50 mil generalistas de 2001 a 2013: \num{a proximidade do
      lugar de nascimento ou de formação é o principal fator}; salário e
      infraestrutura pesam menos.

      \vspace{0.4em}
      \textbf{Austrália}, 3.727 clínicos: \num{65\% não mudariam por pacote
      nenhum}.
    \end{block}

    \vspace{1.0em}
    {\small \lead{Nenhuma dessas desvantagens se resolve sozinha.} Em 2025 o
    Ministério da Saúde reconheceu \num{situação de urgência em saúde pública},
    por \textbf{24 meses}, em razão do tempo de espera na atenção especializada
    --- e lançou o \textbf{Agora Tem Especialistas}, do qual o \textbf{PMM-E} é
    o braço de provimento.}
  }%

  \fonte{\textbf{Fontes:} Scheffer et al., \textit{Demografia Médica no Brasil
  2025} (FMUSP/AMB), dez/2024; Portaria GM/MS nº 7.061/2025; Costa, Nunes \&
  Sanches (2024), \textit{REStat}; Scott et al. (2013), \textit{Social Science
  \& Medicine}.}
\end{frame}

% ===========================================================================
% 4 — O que é o PMM-E                                        (1 frame)
% ===========================================================================
\subsection{A política}
\rastreio{1. Motivação \textperiodcentered\ A política \textperiodcentered\ 1 de 2}

\begin{frame}{O que é o PMM-E}
  \only<1>{%
    \begin{block}{Lei}
      \small
      A \textbf{Lei nº 15.233/2025} criou o \textbf{Projeto Mais Médicos
      Especialistas} dentro do Programa Mais Médicos, \textit{``destinado ao
      provimento de profissionais com vistas à redução no tempo de espera''} do
      usuário do SUS \textit{``nas regiões prioritárias''}.
    \end{block}

    \begin{block}{Quem}
      \small
      Médicos com diploma brasileiro ou revalidado e \textbf{RQE} na área da
      vaga. Não é concurso nem emprego: recebem \num{bolsa-formação} mensal do
      Ministério, \textbf{sem vínculo}.
    \end{block}

    \begin{block}{O quê}
      \small
      \textbf{Aprimoramento em serviço} de \num{12 meses}, \num{20 horas
      semanais} em estabelecimento do SUS, com supervisão e mentoria de
      instituição formadora. São \textbf{16 cursos} --- 6 cirúrgicos e 10
      ambulatoriais --- com foco no câncer e no diagnóstico que o SUS mais
      espera.
    \end{block}
  }%
  \only<2>{%
    \vspace{0.3em}
    \begin{block}{Onde, no primeiro ciclo (julho de 2025)}
      \small
      \num{1.295 vagas} estabelecimento--curso em \num{460 estabelecimentos} e
      \num{368 municípios}, em todas as UFs. O \textbf{Nordeste} concentra 39\%
      das vagas; dois terços dos municípios têm menos de 100 mil habitantes;
      \textbf{18 são capitais}.
    \end{block}

    \vspace{1.2em}
    \begin{destaque}[0.93\textwidth]
      O Ministério publica o quadro de vagas --- município, estabelecimento e
      curso --- e cada vaga sai com uma \textbf{faixa de bolsa}.\\[0.35em]
      É essa regra --- e só ela --- que o trabalho estuda.
    \end{destaque}
  }%

  \fonte{\textbf{Fontes:} Lei nº 15.233/2025, art. 22-D; Portaria GM/MS
  nº 7.177/2025; Edital SGTES/MS nº 3/2025 (DOU 24/07/2025), itens 1, 3--5, 10
  e 11 e Tabela 3; quadro de vagas do ciclo 1, chamada 1; Censo 2022 (IBGE).}
\end{frame}

% ===========================================================================
% 5 — A bolsa remunera o lugar                               (1 frame)
% ===========================================================================
\rastreio{1. Motivação \textperiodcentered\ A política \textperiodcentered\ 2 de 2}

\begin{frame}{A bolsa remunera o lugar}
  \only<1>{%
    \vspace{0.2em}
    {\small O valor da bolsa \textbf{não} depende da especialidade, da carga nem
    do que o médico produz. Depende de \lead{onde fica o município}, em três
    passos:}

    \vspace{0.7em}
    {\small
    \begin{enumerate}
      \setlength{\itemsep}{0.55em}
      \item \textbf{O índice.} O \textbf{IVS 2010 do Ipea} resume, em um número
        de 0 a 1, dezesseis indicadores do Censo 2010 em três dimensões:
        \textit{infraestrutura urbana} (saneamento, lixo, tempo de deslocamento),
        \textit{capital humano} (mortalidade infantil, analfabetismo, crianças
        fora da escola) e \textit{renda e trabalho} (pobreza, desemprego,
        informalidade).
      \item \textbf{A categoria.} O Ipea corta o índice em cinco categorias:
        muito baixa (até 0,200), baixa (0,201--0,300), média (0,301--0,400),
        alta (0,401--0,500) e muito alta (acima de 0,500).
      \item \textbf{O valor.} O edital de 2025 agrupou as cinco categorias em
        \textbf{três faixas}.
    \end{enumerate}}
  }%
  \only<2>{%
    \begin{columns}[c]
      \begin{column}{0.44\textwidth}
        {\footnotesize\renewcommand{\arraystretch}{0.95}
        \begin{tabular}{@{}l c r@{}}
          \toprule
          \makecell[l]{\bfseries Categoria\\\bfseries de IVS} &
          \bfseries Faixa &
          \makecell[r]{\bfseries Bolsa\\\bfseries mensal} \\
          \midrule
          muito alta & 1 & R\$ 20.000 \\
          alta       & 2 & R\$ 15.000 \\
          \makecell[l]{média, baixa\\ou muito baixa} & 3 & R\$ 10.000 \\
          \bottomrule
        \end{tabular}

        No ciclo 1: \num{102} municípios na Faixa 1, \num{107} na Faixa 2 e
        \num{159} na Faixa 3.}
      \end{column}
      \begin{column}{0.53\textwidth}
        \includegraphics[width=\textwidth,height=0.34\textheight,keepaspectratio]%
          {bolsa_por_faixa}
      \end{column}
    \end{columns}

    {\small\lead{O IVS é o piso, não o critério.}}

    \begin{columns}[T]
      \begin{column}{0.49\textwidth}
        \begin{beamercolorbox}[wd=\textwidth,sep=0.3em,rounded=true]{block body}
          \footnotesize \textbf{Duas cláusulas.} A 11.1.4 dá a tabela. A
          \textbf{11.1.3} manda seguir ``critérios de \emph{localização}'' do
          \textbf{Anexo IV}, que o edital não reproduz.
        \end{beamercolorbox}
      \end{column}
      \begin{column}{0.49\textwidth}
        \begin{beamercolorbox}[wd=\textwidth,sep=0.3em,rounded=true]{block body}
          \footnotesize \textbf{Uma direção só.} Em \num{177 dos 368} a faixa
          publicada não é a da categoria: \textbf{zero} abaixo, \textbf{177}
          acima. A categoria é piso; a localização promove 48\%.
        \end{beamercolorbox}
      \end{column}
    \end{columns}
  }%

  \fonte{\textbf{Fontes:} Edital SGTES/MS nº 3/2025, itens 11.1.3 e 11.1.4, e
  retificação; Chamamento SGTES/MS nº 1/2026; Ipea, \textit{Atlas da
  Vulnerabilidade Social nos Municípios Brasileiros} (2015); quadro de vagas do
  ciclo 1. Vale a faixa publicada.}
\end{frame}

% ===========================================================================
% 6 — Onde a bolsa é maior, o médico fica sozinho              (1 frame)
% As duas figuras sustentam a mesma afirmacao; ficam em overlay para ter a
% largura inteira do frame.
% ===========================================================================
\subsection{O efeito é incerto}
\rastreio{1. Motivação \textperiodcentered\ O efeito é incerto \textperiodcentered\ 1 de 3}

\begin{frame}{Onde a bolsa é maior, o médico fica sozinho}
  \vspace{0.2em}
  \only<1>{%
    {\small \lead{Onde a regra manda o dinheiro?} Pela faixa \emph{publicada},
    nos \textbf{295 municípios} com vaga: \lead{por habitante, a bolsa maior
    não vai para onde falta mais.}}

    \vspace{0.3em}
    \centerline{%
      \includegraphics[width=\textwidth,height=0.50\textheight,keepaspectratio]%
        {oferta_pre_por_faixa}}
  }%
  \only<2>{%
    {\small \lead{Em número de colegas, vai.} Mediana de \textbf{2,5} colegas
    da mesma especialidade na Faixa 1, contra \textbf{6,5} na Faixa 3. \lead{O
    programa põe R\$ 5 mil a mais onde o médico trabalharia sozinho.}}

    \vspace{0.3em}
    \centerline{%
      \includegraphics[width=\textwidth,height=0.50\textheight,keepaspectratio]%
        {retaguarda_por_faixa}}
  }%

  \fonte{\textbf{Fontes:} CNES 06/2025 --- profissionais nos CBOs dos 10 cursos
  com correspondência unívoca curso--CBO, nos 295 municípios com vaga no ciclo
  1; população do Censo 2022 (IBGE); \textbf{faixa pela bolsa publicada}, não
  pela categoria de IVS recalculada --- divergem em 177 dos 368 municípios.}
\end{frame}

% ===========================================================================
% 7 — A evidência não decide se R$ 5 mil bastam               (1 frame)
% ===========================================================================
\rastreio{1. Motivação \textperiodcentered\ O efeito é incerto \textperiodcentered\ 2 de 3}

\begin{frame}{A evidência não decide se R\$ 5 mil bastam}
  \only<1>{%
    {\footnotesize O programa aposta que dinheiro compensa lugar ruim.
    \lead{A aposta tem precedente --- e tem limites.}}

    \vspace{-1.0em}
    \begin{columns}[T]
      \begin{column}{0.485\textwidth}
        \begin{block}{A favor: pagar mais funciona}
          \footnotesize\linespread{0.95}\selectfont
          \textbf{No México, o salário foi sorteado.} Um concurso público real
          distribuiu \textbf{106 postos} em municípios pobres e anunciou,
          \textbf{ao acaso}, dois salários. Onde era \textbf{33\% maior}, a
          aceitação subiu \num{15 pontos percentuais}; a mais de \textbf{200 km}
          da cidade natal, foi de \num{25\% para cerca de 80\%} --- \textbf{sem}
          atrair candidatos menos qualificados ou motivados.

          \textbf{O degrau do PMM-E tem o tamanho que a literatura pede.} O que
          um médico exige para ir a um posto pior vai de \num{37\% a 64\%} da
          renda anual em cidades pequenas (Austrália); a oferta no interior
          brasileiro responde a salário com elasticidade de \num{0,7}. O degrau
          --- R\$ 5 mil sobre R\$ 10 mil, ou \textbf{+50\%} --- cai nessa faixa.
        \end{block}
      \end{column}
      \begin{column}{0.485\textwidth}
        \begin{block}{Contra: é caro, e não segura}
          \footnotesize\linespread{0.95}\selectfont
          \textbf{Muitos não vão por preço nenhum.} Dos \textbf{3.727 clínicos
          australianos}, \num{65\%} ficaram onde estavam em todos os cenários.
          Para o pior posto, quem mudaria pedia \num{130\%} da renda anual.

          \vspace{0.2em}
          \textbf{No Brasil, salário compra pouco e custa muito.} Aumentar em
          \textbf{50\%} o salário público no interior do Norte e do Nordeste
          corrigiria \num{12,4\%} do desequilíbrio, a \num{US\$ 15,7 milhões por
          ponto percentual}. Reservar vagas nas faculdades para quem nasceu
          nessas regiões corrigiria \num{63,8\%}, por \num{US\$ 2,2 a 5,1
          milhões} o ponto.
        \end{block}
      \end{column}
    \end{columns}
  }%
  \only<2>{%
    \begin{block}{E o médico vai embora quando a obrigação acaba}
      \small
      Nos \textbf{Estados Unidos}, oito anos depois, \num{12\%} dos médicos que
      foram para clínicas rurais com bolsa e \textbf{obrigação de permanência}
      ainda estavam lá --- contra \num{39\%} dos que foram \textbf{sem
      obrigação nenhuma}.
    \end{block}

    {\small\textbf{Ressalva:} os percentuais da literatura são sobre a
    \textbf{renda total} do médico; a bolsa do PMM-E remunera \textbf{20 horas
    semanais}.}

    \vspace{-0.3em}
    \begin{destaque}[0.93\textwidth]
      Dinheiro move alocação, mas é caro, não move todo mundo e não garante
      que quem foi fique.\\[0.35em]
      A evidência \textbf{não decide} se um degrau de R\$ 5 mil basta.
    \end{destaque}
  }%

  \fonte{\textbf{Fontes:} Dal Bó, Finan \& Rossi (2013), \textit{QJE}; Scott
  et al. (2013), \textit{Social Science \& Medicine}; Costa, Nunes \& Sanches
  (2024), \textit{REStat}; Pathman, Konrad \& Ricketts (1992), \textit{JAMA}.}
\end{frame}

% ===========================================================================
% 8 — No primeiro ciclo, a bolsa maior não ordenou...        (1 frame)
% ===========================================================================
\rastreio{1. Motivação \textperiodcentered\ O efeito é incerto \textperiodcentered\ 3 de 3}

\begin{frame}{No primeiro ciclo, a bolsa maior não ordenou o preenchimento}
  {\small Das \num{1.295 vagas} da primeira chamada, \num{30\%} tiveram alguém
  confirmado ou homologado.}

  \begin{columns}[T]
    \begin{column}{0.56\textwidth}
      \includegraphics[width=\textwidth,height=0.43\textheight,keepaspectratio]%
        {preenchimento_ciclo1}
    \end{column}
    \begin{column}{0.41\textwidth}
      {\small
      \begin{itemize}
        \setlength{\itemsep}{0.15em}
        \item Por \textbf{faixa de bolsa}: \lead{pagar o dobro não preencheu
          mais que pagar uma vez e meia.}
        \item Por \textbf{território}: o preenchimento cai da capital e da
          região metropolitana para o interior ligado a um polo, e é
          \lead{menor no interior remoto.}
      \end{itemize}}

      \vspace{0.1em}
      \begin{beamercolorbox}[wd=\textwidth,sep=0.3em,rounded=true]{block body}
        \footnotesize Isso é \textbf{descrição, não efeito}: as faixas diferem
        em muito mais do que na bolsa. Mas basta para colocar a pergunta.
      \end{beamercolorbox}
    \end{column}
  \end{columns}

  \fonte{\textbf{Fontes:} quadro de vagas e resultados do ciclo 1, chamada 1
  (Ministério da Saúde, 2025); estratos pela REGIC 2018 e pelas RMs e RIDEs
  de 2022 (IBGE).}
\end{frame}

% ===========================================================================
% 9 — Pergunta                                               (1 frame)
% ===========================================================================
\section{Pergunta}
\rastreio{2. Pergunta}

\begin{frame}{Pergunta}
  \vspace{0.5em}
  \begin{destaque}
    \large Maiores bolsas do PMM-E para municípios mais vulneráveis
    compensam suas desvantagens territoriais na atração de médicos
    especialistas?
  \end{destaque}

  \vspace{0.8em}
  {\small Dois objetos, um contra o outro:}

  \vspace{0.4em}
  \begin{columns}[T]
    \begin{column}{0.47\textwidth}
      \begin{beamercolorbox}[wd=\textwidth,sep=0.75em,rounded=true]{block body}
        \begin{minipage}[t][3.0em][c]{\dimexpr\textwidth-1.5em\relax}
          \small\textbf{O preço} que a política colocou sobre a
          vulnerabilidade --- o degrau de \textbf{R\$ 5 mil} entre faixas
        \end{minipage}
      \end{beamercolorbox}
    \end{column}
    \begin{column}{0.47\textwidth}
      \begin{beamercolorbox}[wd=\textwidth,sep=0.75em,rounded=true]{block body}
        \begin{minipage}[t][3.0em][c]{\dimexpr\textwidth-1.5em\relax}
          \small\textbf{A desvantagem} que esse preço pretende compensar
        \end{minipage}
      \end{beamercolorbox}
    \end{column}
  \end{columns}

  \vspace{0.7em}
  {\small A margem que observamos é o \lead{preenchimento da vaga}: se, ao
  final da chamada, apareceu alguém disposto a ocupá-la.}
\end{frame}

% ===========================================================================
% 10 — A decisão: onde vale a pena estar                     (1 frame)
% ===========================================================================
\section{Literatura teórica}
\rastreio{3. Literatura teórica \textperiodcentered\ 1 de 2}

\begin{frame}{A decisão: onde vale a pena estar}
  \vspace{0.2em}
  {\small \textbf{Moehling, Niemesh, Thomasson \& Treber (2020)}, eq. 1,
  p. 184, dão a \lead{estrutura da decisão}: escolhe-se a localidade que
  maximiza o valor presente do rendimento \textbf{real}, líquido do custo não
  pecuniário de viver ali.}

  \vspace{-0.6em}
  \[
    \arg\max_{i \in I}\ \left\{ \sum_t \delta^t
    \left[ \dfrac{\mathbb{E}\!\left(w_{it}^{(s)}\right)}{p_{it}}
    - c_{it}^{(s)} \right] \right\}
  \]
  \vspace{-1.3em}

  {\footnotesize\renewcommand{\arraystretch}{1.1}
  \begin{tabularx}{\textwidth}{@{}>{\centering\arraybackslash}p{0.19\textwidth} L@{}}
    \toprule
    \bfseries Termo & \bfseries Leitura \\
    \midrule
    $\sum_t \delta^t$ & a escolha é de \textbf{carreira}: fixar-se ou migrar ao
      fim do vínculo \\
    $\mathbb{E}(w^{(s)}_{it}) / p_{it}$ & rendimento \textbf{deflacionado} pelo
      custo de vida local \\
    $c^{(s)}_{it}$ & o que torna estar ali custoso e \textbf{não é pago em
      dinheiro} \\
    $s$ & generalista atende em posto simples; \textbf{especialista} precisa de
      centro cirúrgico e leito \\
    \bottomrule
  \end{tabularx}}

  \vspace{0.3em}
  {\footnotesize Na definição dos próprios autores, $c$ reúne
  \textit{``preferences over rural or urban living, or other location-specific
  attributes, such as proximity to family''} --- uma \lead{caixa-preta}. As
  duas referências seguintes a abrem.}
\end{frame}

% ===========================================================================
% 11 — Abrindo o custo: o lugar e o trabalho                 (1 frame)
% ===========================================================================
\rastreio{3. Literatura teórica \textperiodcentered\ 2 de 2}

\begin{frame}{Abrindo o custo: o lugar e o trabalho}
  \only<1>{%
    \vspace{0.2em}
    {\small \lead{O lugar.} \textbf{Redding \& Rossi-Hansberg (2017)}, eq. 24,
    p. 28, dão o \textbf{custo geográfico}: a utilidade de trabalhar num lugar
    depende do salário, das \textbf{amenidades} e do \textbf{custo de
    moradia}.}

    \vspace{-0.4em}
    \[ u_{nio} = \dfrac{z_{nio}\, B_n\, w_i}{\kappa_{ni}\, Q_n^{\,1-\beta}}
       \quad\Longrightarrow\quad
       c^{\text{espacial}}_m = (1-\beta)\ln Q_m - \ln A_m \]
    \vspace{-1.0em}

    {\small Somado à proximidade da família de Moehling et al.:}

    \vspace{-0.5em}
    \[ c^{\text{geo}}_{im} = \phi(\text{dist}_{im}) - \gamma A_m
       + \theta_i^{\text{rural}} \]
    \vspace{-1.0em}

    {\footnotesize\renewcommand{\arraystretch}{1.2}
    \begin{tabularx}{\textwidth}{@{}>{\centering\arraybackslash}p{0.16\textwidth} >{\centering\arraybackslash}p{0.11\textwidth} L@{}}
      \toprule
      \bfseries Componente & \bfseries Sinal & \bfseries Significado \\
      \midrule
      $\phi(\text{dist}_{im})$ & $\phi' > 0$ & afastar-se da família custa, e
        custa mais a cada quilômetro \\
      $-\gamma A_m$ & $< 0$ & amenidade urbana --- saneamento, segurança, escola
        --- compensa \\
      $\theta_i^{\text{rural}}$ & $\gtrless 0$ & gosto por cidade pequena ou
        grande, sem sinal universal \\
      \bottomrule
    \end{tabularx}}

    \vspace{0.5em}
    {\footnotesize \textbf{Limitação assumida.} CNES e edital não informam a
    residência do profissional, por sigilo fiscal. A unidade é o
    \textbf{município do estabelecimento}: $\text{dist}$ entra como latente.}
  }%
  \only<2>{%
    \vspace{0.2em}
    {\small \lead{O trabalho.} \textbf{Choné \& Ma (2011)}, eq. 1, p. 232, dão o
    \textbf{custo laboral}: soma-se a renda, subtrai-se o custo de atender e
    soma-se o benefício ao paciente, ponderado pelo \textbf{altruísmo}.}

    \vspace{-0.6em}
    \[ U = R - C(q; L, K) + \alpha B(q)
       \quad\Longrightarrow\quad
       c^{\text{laboral}}_{im} = C(q; L_m, K_m) - \alpha_i B(q; L_m, K_m) \]
    \vspace{-1.3em}

    {\footnotesize\renewcommand{\arraystretch}{1.05}
    \begin{tabularx}{\textwidth}{@{}>{\centering\arraybackslash}p{0.14\textwidth} >{\centering\arraybackslash}p{0.22\textwidth} L@{}}
      \toprule
      \bfseries Componente & \bfseries Derivadas & \bfseries Significado \\
      \midrule
      $C(q)$ & $C' > 0,\ C'' > 0$ & atender cansa, e cansa de forma
        \textbf{crescente} \\
      $\alpha_i B(q)$ & $B' > 0,\ B'' < 0$ & curar dá satisfação, decrescente
        porque a triagem prioriza o caso grave \\
      \bottomrule
    \end{tabularx}}

    \vspace{0.3em}
    {\footnotesize O custo marginal $c'(q) = C' - \alpha B'$ tem \textbf{sinal
    incerto}, mas $c'' \gg 0$: a curva é um \textbf{U} --- uma zona em que
    atender mais \emph{reduz} o custo líquido, um mínimo, e uma zona de
    exaustão.}

    \vspace{0.3em}
    {\footnotesize \lead{Equipe ($L$) e capital ($K$) atuam duas vezes.}
    Reduzem o cansaço, $\partial C/\partial K < 0$ --- canal que já está em
    Choné \& Ma. E \textbf{ampliam o benefício}, $\partial B/\partial K > 0$:
    sem medicamento, insumo cirúrgico ou maquinário em funcionamento, o
    atendimento perde resolutividade --- \textbf{extensão deste projeto},
    motivada por Reinhardt (1972, 1975). Por dois caminhos,
    $\partial c^{\text{laboral}}/\partial K < 0$.}
  }%
\end{frame}

% ===========================================================================
% 12 — Como juntamos os três                                 (1 frame)
% ===========================================================================
\section{Modelo microeconômico}
\rastreio{4. Modelo microeconômico \textperiodcentered\ 1 de 2}

\begin{frame}{Como juntamos os três}
  \only<1>{%
    \vspace{0.2em}
    {\small A estrutura vem de Moehling et al.; o custo, que neles era
    caixa-preta, é aberto pelas outras duas.}

    \vspace{-0.3em}
    \[ V_{im} = \sum_t \delta^t \left[ \frac{\mathbb{E}(w_{imt} \mid B_m)}{p_{mt}}
       - c_{im} \right] + \varepsilon_{im} \]
    \vspace{-1.4em}
    \[ c_{im} = \underbrace{\phi(\text{dist}_{im}) - \gamma A_m
       + \theta_i^{\text{rural}}}_{\text{Redding \& Rossi-Hansberg}}
       + \underbrace{C(q; L_m, K_m)
       - \alpha_i B(q; L_m, K_m)}_{\text{Choné \& Ma, com a extensão em } B} \]
    \vspace{-1.0em}

    {\footnotesize\renewcommand{\arraystretch}{1.25}
    \begin{tabularx}{\textwidth}{@{}>{\raggedright\arraybackslash}p{0.32\textwidth} L@{}}
      \toprule
      \bfseries De onde vem & \bfseries O que entrega \\
      \midrule
      Moehling et al. (2020) & o $\arg\max$ intertemporal, o deflator $p_{mt}$
        e a existência de $c$ \\
      Redding \& Rossi-Hansberg (2017) & distância, amenidades e custo de
        moradia dentro de $c$ \\
      Choné \& Ma (2011), com Reinhardt & esforço, altruísmo e o papel de $L$ e
        $K$ dentro de $c$ \\
      \bottomrule
    \end{tabularx}}
  }%
  \only<2>{%
    \vspace{1.0em}
    \begin{beamercolorbox}[wd=\textwidth,sep=0.9em,rounded=true]{block body}
      \small \textbf{A decisão.} A alternativa $m = 0$ é \textbf{ficar fora do
      programa}. A vaga em $m$ é aceita quando $V_{im} \geq V_{i0}$ e $m$ é a
      melhor entre as disponíveis. O termo $\varepsilon_{im}$ recolhe os gostos
      que não observamos.
    \end{beamercolorbox}

    \vspace{1.6em}
    \begin{destaque}[0.93\textwidth]
      Nenhuma das três trata de um componente da remuneração fixado por
      \textbf{regra pública sobre um índice territorial}. É isso, e só isso,
      que a adaptação ao PMM-E acrescenta.
    \end{destaque}

    \vspace{1.2em}
    {\centering\small É o que o slide seguinte desenvolve.\par}
  }%
\end{frame}

% ===========================================================================
% 13 — O IVS organiza o custo                                (1 frame)
% ===========================================================================
\rastreio{4. Modelo microeconômico \textperiodcentered\ 2 de 2}

\begin{frame}{O IVS organiza o custo}
  \only<1>{%
    {\small \lead{O que a bolsa paga --- e o que não paga.} A remuneração tem
    duas partes: a \textbf{bolsa}, que a regra fixa, e o que o médico obtém no
    \textbf{mercado local} fora das 20 horas do programa:}

    \vspace{-0.9em}
    \[ \mathbb{E}(w_{imt} \mid B_m) = B_m + w^{\text{priv}}_m \]
    \vspace{-1.4em}

    \begin{columns}[T]
      \begin{column}{0.49\textwidth}
        \begin{beamercolorbox}[wd=\textwidth,sep=0.4em,rounded=true]{block body}
          \footnotesize \textbf{Na capital}, $w^{\text{priv}}$ é alto: R\$ 10
          mil mais o consultório pode superar R\$ 20 mil sem complemento.
        \end{beamercolorbox}
      \end{column}
      \begin{column}{0.49\textwidth}
        \begin{beamercolorbox}[wd=\textwidth,sep=0.4em,rounded=true]{block body}
          \footnotesize \textbf{No interior isolado}, $w \to B$: a bolsa é
          \textbf{toda} a remuneração --- e é ali que a política mais aposta
          nela.
        \end{beamercolorbox}
      \end{column}
    \end{columns}

    \vspace{0.3em}
    {\small O deflator $p_{mt}$ trabalha no sentido oposto: o custo de vida é
    \textbf{menor} no interior, o que valoriza a mesma bolsa em termos reais.}
  }%
  \only<2>{%
    {\small \lead{O que se observa.} Distância da família, aluguel, mercado
    privado e esforço clínico \textbf{não são observados}. O que se observa,
    para todo município, é o \textbf{IVS}. Escrevemos o custo como função do
    índice mais um desvio individual:}

    \vspace{-0.9em}
    \[ c_{im} = c_0(IVS_m) + \eta_i \]
    \vspace{-1.4em}

    {\small Isso não é atalho: \lead{cada dimensão do IVS corresponde a um
    bloco do custo.}}

    \vspace{0.2em}
    {\small\renewcommand{\arraystretch}{1.1}
    \begin{tabularx}{\textwidth}{@{}>{\raggedright\arraybackslash}p{0.16\textwidth} >{\raggedright\arraybackslash}p{0.24\textwidth} L >{\centering\arraybackslash}p{0.11\textwidth}@{}}
      \toprule
      \bfseries Dimensão & \bfseries Indicadores & \bfseries Bloco do custo &
        \bfseries Efeito \\
      \midrule
      Infraestrutura urbana & saneamento, lixo, tempo de deslocamento &
        amenidades $A_m$ & $\uparrow$ \\
      Renda e trabalho & pobreza, desemprego, informalidade & mercado privado
        ausente, $w \to B$ & $\uparrow$ \\
      Capital humano & mortalidade infantil, analfabetismo, mães adolescentes &
        gravidade do caso, $B'(q)\uparrow$; escassez de equipe e capital,
        $K\downarrow$ & \textbf{ambíguo} \\
      \bottomrule
    \end{tabularx}}
  }%
  \only<3>{%
    \vspace{0.2em}
    {\small\lead{A terceira linha --- capital humano --- é o que impede assumir
    que o custo cresce com o índice.}}

    \vspace{0.6em}
    \begin{columns}[T]
      \begin{column}{0.47\textwidth}
        \begin{beamercolorbox}[wd=\textwidth,sep=0.6em,rounded=true]{block body}
          \begin{minipage}[t][3.4em][c]{\dimexpr\textwidth-1.7em\relax}
            \small Carência sanitária \textbf{eleva o benefício} de atender ---
            o que \textbf{reduz} o custo para um médico altruísta.
          \end{minipage}
        \end{beamercolorbox}
      \end{column}
      \begin{column}{0.47\textwidth}
        \begin{beamercolorbox}[wd=\textwidth,sep=0.6em,rounded=true]{block body}
          \begin{minipage}[t][3.4em][c]{\dimexpr\textwidth-1.7em\relax}
            \small Ao mesmo tempo \textbf{sinaliza falta de insumos} --- o que
            \textbf{eleva} o cansaço.
          \end{minipage}
        \end{beamercolorbox}
      \end{column}
    \end{columns}

    \vspace{0.6em}
    {\centering\large $c_0'(IVS) \gtrless 0$ \quad {\normalsize o sinal é
    questão empírica}\par}

    \vspace{0.7em}
    \begin{destaque}[0.93\textwidth]
      É por isso que o objeto do trabalho é o \textbf{degrau} da bolsa entre
      faixas, e não a inclinação do índice.
    \end{destaque}
  }%
\end{frame}

% ===========================================================================
% 14 — A hipótese                                            (1 frame)
% ===========================================================================
\section{Hipótese}
\rastreio{5. Hipótese}

\begin{frame}{A hipótese}
  \only<1>{%
    \vspace{0.5em}
    {\small\lead{Passo 1 --- a condição de aceitação.} O médico $i$ aceita a
    vaga em $m$ quando o que ela vale supera sua melhor alternativa
    $\bar{v}_i$:}

    \vspace{-0.4em}
    \[
      \frac{B_m + w^{\text{priv}}_m}{p_m} - c_0(IVS_m) \geq \bar{v}_i
    \]
    \vspace{-0.9em}

    {\small\lead{Passo 2 --- do médico à vaga.} A vaga é preenchida se existir
    \textbf{ao menos um} candidato para quem a desigualdade vale. Tudo o que
    aumenta o lado esquerdo aumenta essa probabilidade.}

    \vspace{0.4em}

    \vspace{1.8em}
    {\small\lead{Passo 3 --- a hipótese.} O trabalho testa uma só.}

    \vspace{0.8em}
    {\small
    \begin{tabularx}{\textwidth}{@{}>{\centering\arraybackslash}p{0.06\textwidth} L >{\centering\arraybackslash}p{0.32\textwidth}@{}}
      \toprule
      & \bfseries Hipótese & \bfseries Derivada \\
      \midrule
      \textbf{H1} & Maior \textbf{remuneração real} aumenta a probabilidade de
        preenchimento da vaga
        & $\dfrac{\partial \Pr(\text{preench.}_m)}{\partial (B_m / p_m)} > 0$ \\
      \bottomrule
    \end{tabularx}}
  }%
  \only<2>{%
    \vspace{0.8em}
    {\small\lead{Passo 4 --- o custo é obstáculo, não hipótese.} O custo
    locacional não entra como segunda hipótese porque não varia livremente: a
    regra do edital o amarra à bolsa, e os dois sobem juntos. Ele é o que torna
    \textbf{H1} difícil de testar, não uma segunda afirmação a testar.}

    \vspace{1.3em}
    {\small Na fronteira entre duas faixas, o preenchimento do lado mais
    vulnerável exige que o \textbf{degrau monetário} supere o \textbf{degrau de
    custo}:}

    \vspace{1.0em}
    \begin{center}
      \begin{minipage}{0.66\textwidth}
        \begin{beamercolorbox}[wd=\textwidth,sep=0.9em,center,rounded=true]{block body}
          $\displaystyle \frac{\Delta B_m}{p_m} > \Delta c_0,
           \quad \Delta B_m = \text{R\$ }5.000$
        \end{beamercolorbox}
      \end{minipage}

      \vspace{1.4em}
      {\large A pergunta da apresentação é se essa desigualdade vale.}
    \end{center}
  }%
\end{frame}

% ===========================================================================
% 15 — Viabilidade empírica                                  (1 frame)
% ===========================================================================
\section{Viabilidade empírica}
\rastreio{6. Viabilidade empírica}

\begin{frame}{Viabilidade empírica}
  \only<1>{%
    {\small\lead{Temos como medir cada peça do modelo?}}

    \vspace{0.2em}
    {\footnotesize\renewcommand{\arraystretch}{1.0}
    \begin{tabularx}{\textwidth}{@{}>{\raggedright\arraybackslash}p{0.19\textwidth} L >{\raggedright\arraybackslash}p{0.27\textwidth}@{}}
      \toprule
      \bfseries Peça do modelo & \bfseries O que observamos & \bfseries Fonte \\
      \midrule
      Preenchimento da vaga & se teve alguém confirmado ou homologado &
        quadro de vagas e resultados do ciclo 1: 1.295 vagas, 368 municípios \\
      Bolsa $B_m$ & o valor anunciado na vaga: R\$ 10, 15 ou 20 mil & edital \\
      Custo $c_m$ & o IVS e suas três dimensões; capital, metropolitano, polo
        do interior ou interior remoto; especialistas e estrutura
        preexistentes & Ipea; REGIC 2018 e RMs 2022 (IBGE); CNES mensal \\
      Custo de vida $p_m$ & diferenças entre estados & IBGE \\
      Mercado local $w^{\text{priv}}_m$ & --- & \textbf{não observado} \\
      Distância da família & --- & \textbf{não observado} \\
      \bottomrule
    \end{tabularx}}

    \vspace{0.5em}
    {\small\lead{Sim, para o essencial.} As duas peças não observadas --- a
    renda no mercado privado local e a distância da família --- entram no
    modelo como parte do custo que o IVS e a tipologia territorial resumem.}
  }%
  \only<2>{%
    \vspace{0.3em}
    {\small\lead{A dificuldade, e o que já sabemos dela.} O que segue é sobre
    \textbf{o efeito do valor da bolsa} --- o objeto de H1 ---, não sobre o
    trabalho inteiro. Recuperar a regra que o Ministério aplicou era o primeiro
    passo do trabalho empírico. \textbf{Ele foi dado}, e tem três resultados.}

    \vspace{0.5em}
    {\small
    \begin{enumerate}
      \setlength{\itemsep}{0.5em}
      \item \textbf{Há suporte comum.} \num{37} municípios com IVS $\leq$ 0,400
        estão na Faixa 1 e \num{94} na Faixa 2. Os intervalos de IVS das três
        faixas se sobrepõem. Não falta variação.
      \item \textbf{Mas não há descontinuidade onde os cortes estão.} Em 0,500
        os dois lados são \textbf{100\% Faixa 1} em qualquer janela até
        $\pm$0,050; em 0,400 não há Faixa 3 por perto --- o maior IVS da Faixa 3
        é 0,372. O tratamento é \textbf{localmente constante nos dois cortes}.
      \item \textbf{E a variação que sobra não é exógena.} Dos 83 municípios
        fora da melhor regra de limiar, os 41 promovidos têm mediana de
        população de \num{7.933} contra \num{32.179} dos 42 rebaixados, e muito
        mais interior remoto. Parear por IVS confunde bolsa com posição
        territorial --- contra a bolsa.
    \end{enumerate}}
  }%
  \only<3>{%
    \vspace{0.8em}
    \begin{destaque}[0.93\textwidth]
      O efeito causal do \emph{valor da bolsa} não é identificável com as
      fontes públicas; a razão não é potência amostral. Destravá-lo exige o
      \textbf{Anexo IV} e os critérios de localização da cláusula 11.1.3, que
      só o Ministério pode fornecer.
    \end{destaque}

    \vspace{1.2em}
    \begin{block}{O que fica de pé}
      \small
      A pergunta do slide 9 continua sendo a pergunta do trabalho. Até que a
      regra seja recuperada, a relação entre faixa e preenchimento é reportada
      como \textbf{gradiente}, não como efeito.

      \vspace{0.4em}
      O trabalho tem um desenho causal que não depende disso --- a
      descontinuidade no \textbf{escore de seleção do candidato} ---, mas ele
      responde a outra pergunta: o efeito de \textbf{ganhar a vaga}, não o do
      valor da bolsa.
    \end{block}
  }%
\end{frame}

\end{document}
